% Hella Programming Language Specification
% Source grammar: hella/references/ebnf-0.1.txt (Draft 0.4: EBNF 0.1
% plus the async extension: `async' functions, `task<T>', `await',
% `spawn', `scope do ... end', `yield').
\documentclass[makeidx]{article}
\usepackage{xspace}
\usepackage{epsfig}
\usepackage{xcolor}
\usepackage{syntax}
\usepackage[fleqn]{amsmath}
\usepackage{amssymb}
\usepackage{semantic}
\usepackage{hella}
\usepackage[hidelinks]{hyperref}
\usepackage{lmodern}
\usepackage[T1]{fontenc}
\usepackage{makeidx}
\usepackage{tocloft}
\makeindex
\title{Hella Programming Language Specification\\
{\large Draft 0.4 --- EBNF 0.1 plus async}}
\author{Hassan Qasemi}

\begin{document}
\maketitle
\setlength{\cftsubsecnumwidth}{2.7em}
\setlength{\cftsubsubsecnumwidth}{3.4em}
\tableofcontents
\newpage
\pagestyle{myheadings}
\markright{Hella Programming Language Specification}

\section{Scope}
\LMLabel{ecmaScope}

\LMHash{}%
This standard specifies the syntax and semantics of
the Hella programming language.
It does not specify the APIs of the Hella libraries
except where those library elements are essential to
the correct functioning of the language itself.

\section{Conformance}
\LMLabel{ecmaConformance}

\LMHash{}%
A conforming implementation of the Hella programming language
must provide and support all the APIs
mandated in this specification.

\LMHash{}%
A conforming implementation is permitted to provide additional APIs,
but not additional syntax,
except for experimental features.

\section{Normative References}
\LMLabel{ecmaNormativeReferences}

\LMHash{}%
The following referenced documents are indispensable for
the application of this document.
\begin{enumerate}
\item
  The Unicode Standard, Version 5.0, as amended by Unicode 5.1.0, or successor.
\end{enumerate}

\section{Terms and Definitions}
\LMLabel{ecmaTermsAndDefinitions}

\LMHash{}%
Terms and definitions used in this specification are given in
the body of the specification proper.

\section{Notation}
\LMLabel{notation}

\LMHash{}%
We distinguish between normative and non-normative text.
Normative text defines the rules of Hella.
It is given in this font.
At this time, non-normative text includes:

\begin{itemize}
\item[Rationale]
  Discussion of the motivation for language design decisions appears in italics.
\rationale{%
  Distinguishing normative from non-normative helps clarify
  what part of the text is binding and what part is merely expository.%
}
\item[Commentary]
  Comments such as
  ``\commentary{%
    The careful reader will have noticed
    that the name Hella has five characters%
  }''
  serve to illustrate or clarify the specification,
  but are redundant with the normative text.
\commentary{%
  The difference between commentary and rationale can be subtle.%
}
\rationale{%
  Commentary is more general than rationale,
  and may include illustrative examples or clarifications.%
}
\end{itemize}

\LMHash{}%
Reserved words and built-in identifiers
(\ref{identifierReference})
appear in {\bf bold}.

\commentary{%
  Examples would be \CLASS{} or \MATCH.%
}

\LMHash{}%
Grammar productions are given in a common variant of EBNF.
The left hand side of a production ends with `\lit{::=}'.
On the right hand side, alternation is represented by vertical bars,
and sequencing by spacing.
As in PEGs, alternation gives priority to the left.
Optional elements of a production are suffixed by a question mark
like so: \code{anElephant?}.
Appending a star to an element of a production means
it may be repeated zero or more times.
Appending a plus sign to a production means it occurs one or more times.
Parentheses are used for grouping.
Negation is represented by prefixing an element of a production with a tilde.

\commentary{%
  An example would be:%
}

\begin{grammar}\color{commentaryColor}
<aProduction> ::= <anAlternative>
  \alt <anotherAlternative>
  \alt <oneThing> <after> <another>
  \alt <zeroOrMoreThings>*
  \alt <oneOrMoreThings>+
  \alt <anOptionalThing>?
  \alt (<some> <grouped> <things>)
  \alt \gtilde<notAThing>
  \alt `aTerminal'
  \alt <A\_LEXICAL\_THING>
\end{grammar}

\LMHash{}%
Both syntactic and lexical productions are represented this way.
Lexical productions are distinguished by their names.
The names of lexical productions consist exclusively of
upper case characters and underscores.
As always, within grammatical productions,
whitespace and comments between elements of the production
are implicitly ignored unless stated otherwise.
Punctuation tokens appear in quotes.

\LMHash{}%
Productions are embedded, as much as possible,
in the discussion of the constructs they represent.

\LMHash{}%
A \Index{term} is a syntactic construct.
It may be considered to be a piece of text which is derivable in the grammar,
and it may be considered to be a tree created by such a derivation.
An \Index{immediate subterm} of a given term $t$ is a syntactic construct
which corresponds to an immediate subtree of $t$
considered as a derivation tree.
A \Index{subterm} of a given term $t$ is $t$,
or an immediate subterm of $t$,
or a subterm of an immediate subterm of $t$.

\LMHash{}%
A list \DefineSymbol{x_1, \ldots, x_n} denotes any list of
$n$ elements of the form $x_i, 1 \le i \le n$.
Note that $n$ may be zero, in which case the list is empty.
We use such lists extensively throughout this specification.

\BlindDefineSymbol{j, y_j, x_j}%
\LMHash{}%
Let further $y_j$ be an atomic syntactic entity (like an identifier),
$x_j$ a composite syntactic entity (like an expression or a type),
and \DefineSymbol{E} again a composite syntactic entity.

\section{Overview}
\LMLabel{overview}

\LMHash{}%
Hella is a statically typed, class-based language with structs, traits,
enums, extensions, generics with \code{where} constraints, closures,
pattern matching with \code{match}, ranges, nullable and pointer/slice
type modifiers, fixed-size arrays (\code{arr}), growable vectors
(\code{vec}), \code{K:V} maps, owning heap pointers (\code{own T}
with \code{new} construction and \code{delete} destruction), and
structured-concurrency async functions (\code{async} bodies returning
\code{task<T>} handles consumed by \code{await} or produced by
\code{spawn} inside \code{scope} blocks, with \code{yield} checkpoints).
Programs are organized as source files containing top-level declarations
and optional module initialization blocks (\ref{sourceFiles},
\ref{moduleInitialization}).

\LMHash{}%
Hella programs may be statically checked.
Programs with compile-time errors do not have a specified dynamic semantics.
This specification makes no attempt to answer additional questions
about a library or program at the point
where it is known to have a compile-time error.

\commentary{%
  Change log carried over from the source EBNF (Draft 0.4: EBNF 0.1
  Phase 5 fixes plus own-heap phase 1 plus the async extension):
  accessors public by default, separate \code{get}/\code{set} merging,
  \code{open} class-only, omitted \code{Type} in \code{has} literals,
  \code{|}/\code{or} alternative chains plus tuple-pattern \code{(a,b)}
  in \code{match}, \code{arr}/\code{vec} type forms, array literals
  \code{[1,2,3]}, empty \code{vec[]}, \code{K:V} map types plus
  \code{has k: v end} map literals, \code{own T} owning heap pointers
  with \code{new} construction and \code{delete} destruction,
  \code{async} functions with \code{task<T>} handles, prefix
  \code{await}/\code{spawn} expressions, \code{scope} blocks and
  \code{yield} checkpoints (\ref{asyncFunctions}). The \code{=} form
  of attribute arguments and conditional compilation with \code{@cfg}
  (\ref{conditionalCompilation}) are also specified. Extern parameters
  may carry \code{ref}/\code{out} modes for C out-parameters
  (\ref{ffi}). Interpolation in
  strings is written with single braces (not with a dollar prefix),
  and doubled braces denote literal brace characters.%
}

\begin{hellaCode}
\PUBLIC{} \INT{} add(\INT{} a, \INT{} b) \DO{} \RETURN{} a + b; \END{}
\end{hellaCode}

\section{Source Files}
\LMLabel{sourceFiles}

\LMHash{}%
A Hella source file is a sequence of top-level declarations
(\ref{topLevelDeclarations}) and module initialization blocks
(\ref{moduleInitialization}).

\begin{grammar}
<sourceFile> ::= (<topLevelDeclaration> | <initBlock>)*
\end{grammar}

\section{Lexical Grammar}
\LMLabel{lexicalGrammar}

\LMHash{}%
Hella source text is represented as a sequence of Unicode code points.
This sequence is first converted into a sequence of tokens
according to the lexical rules given in this specification.
At any point in the tokenization process,
the longest possible token is recognized.

\begin{grammar}
<identifier> ::= <IDENTIFIER\_START> <IDENTIFIER\_CONTINUE>*

<IDENTIFIER\_START> ::= <LETTER> | `_'

<IDENTIFIER\_CONTINUE> ::= <LETTER> | <DIGIT> | `_'

<LETTER> ::= `a' .. `z' | `A' .. `Z'

<DIGIT> ::= `0' .. `9'
\end{grammar}

\commentary{%
  The EBNF source uses \code{identifier-start} and
  \code{identifier-continue} with a single underscore terminal
  (\code{"\_"}); the lexical rule above derives exactly those forms.%
}

\subsection{Reserved Words}
\LMLabel{reservedWords}

\LMHash{}%
A \Index{reserved word} can only be used in the syntactic positions
specified by the grammar.
In particular, a \Error{compile-time error} occurs if a reserved word is used
where an identifier is expected.
Keywords are reserved and cannot be used as identifiers.

\commentary{%
  Note that reserved words occur bold and unquoted in grammar rules
  (e.g., \CLASS{})
  even though the consistent notation would use quotes
  (e.g., \lit{class}).
  This notational abuse occurs because we believe
  it makes the grammar rules more readable.%
}

\begin{grammar}
<KEYWORD> ::= \gnewline{}
  \AND{} | \ANY{} | \ARR{} | \ASSERT{} | \ASYNC{} | \AWAIT{} | \BOOL{}
    | \BREAK{} | \CLASS{}
  \alt\hspace{-3mm} \CONST{} | \CONTINUE{} | \DEBUGASSERT{} | \DEFER{}
    | \DELETE{} | \DISTINCT{} | \DO{} | \DOUBLE{}
  \alt\hspace{-3mm} \ELSE{} | \ENUM{} | \EXPLICIT{} | \EXTENDS{}
    | \EXTERN{} | \FLOAT{} | \FOR{}
  \alt\hspace{-3mm} \FUNCTION{} | \IF{} | \IMPLEMENTS{} | \IMPORT{}
    | \IN{} | \INIT{} | \INT{}
  \alt\hspace{-3mm} \IS{} | \LOOP{} | \MATCH{} | \NEW{} | \NOT{} | \NULL{}
    | \OPEN{} | \OPERATOR{}
  \alt\hspace{-3mm} \OR{} | \OUT{} | \OWN{} | \OVERRIDE{} | \PRIVATE{}
    | \PUBLIC{} | \REF{} | \RETURN{} | \SCOPE{}
  \alt\hspace{-3mm} \SEALED{} | \SPAWN{} | \STATIC{} | \STRING{} | \STRUCT{}
    | \SUPER{} | \THIS{} | \TRAIT{}
  \alt\hspace{-3mm} \TYPEDEF{} | \VOID{} | \VEC{} | \WHILE{} | \WHERE{}
    | \YIELD{}
\end{grammar}

\LMHash{}%
In the grammar, the rule for reserved words above must occur
before the rule for \synt{identifier}
(\ref{identifierReference}).

\commentary{%
  The EBNF keyword list contains exactly 62 entries:
  \code{and}, \code{any}, \code{arr}, \code{assert}, \code{async},
  \code{await}, \code{bool},
  \code{break}, \code{class}, \code{const}, \code{continue},
  \code{debug\_assert}, \code{defer}, \code{delete}, \code{distinct}, \code{do},
  \code{double}, \code{else}, \code{enum}, \code{extends},
  \code{explicit}, \code{extern}, \code{float}, \code{for},
  \code{function}, \code{if}, \code{implements}, \code{import},
  \code{in}, \code{init}, \code{int}, \code{is}, \code{loop},
  \code{match}, \code{new}, \code{not}, \code{null}, \code{open}, \code{operator},
  \code{or}, \code{out}, \code{own}, \code{override}, \code{private}, \code{public},
  \code{ref}, \code{return}, \code{scope}, \code{sealed}, \code{spawn},
  \code{static}, \code{string},
  \code{struct}, \code{super}, \code{this}, \code{trait}, \code{typedef},
  \code{void}, \code{vec}, \code{while}, \code{where}, \code{yield}.
  The words \code{get}, \code{set}, \code{has}, \code{end},
  \code{convert}, \code{to}, \code{extend}, \code{from},
  \code{initialize} and \code{Self} are contextual, not reserved.
  \code{task} is likewise not reserved: it is only special as the
  \code{task<T>} type form (\ref{types}).%
}

\subsection{Comments and Whitespace}
\LMLabel{comments}

\LMHash{}%
\IndexCustom{Comments}{comment}
are sections of program text that are used for documentation.

\begin{grammar}
<LINE\_COMMENT> ::= \gnewline{}
  `//' \gtilde(<LINE\_BREAK>)* (<LINE\_BREAK>)?

<BLOCK\_COMMENT> ::= `/*' (<ANY\_CHARACTER>)* `*/'

<COMMENT> ::= <LINE\_COMMENT> | <BLOCK\_COMMENT>

<NEWLINE> ::= <LINE\_FEED> | <CARRIAGE\_RETURN\_LINE\_FEED>
  \alt <CARRIAGE\_RETURN>

<WHITESPACE> ::= ` ' | `\\t' | <COMMENT>
\end{grammar}

\LMHash{}%
As always, within grammatical productions,
whitespace and comments between elements of the production
are implicitly ignored unless stated otherwise.
The EBNF non-terminals \code{any-character-except-newline},
\code{any-character}, \code{LF}, \code{CRLF} and \code{CR}
are treated as primitive lexical tokens here.

\section{Literals}
\LMLabel{literals}

\LMHash{}%
A literal denotes a fixed value: a number, a character, a string,
\TRUE{}, \FALSE{} or \NULL{}.

\subsection{Integers}
\LMLabel{integerLiterals}

\LMHash{}%
An \IndexCustom{integer literal}{literal!integer} is either a decimal,
hexadecimal or binary numeral.
Underscores may occur between digits for readability
and carry no numeric meaning.

\begin{grammar}
<integerLiteral> ::= <decimalLiteral>
  \alt <hexadecimalLiteral>
  \alt <binaryLiteral>

<decimalLiteral> ::= <DIGIT> (<DIGIT> | `_')*

<hexadecimalLiteral> ::= `0x' <HEX\_DIGIT> (<HEX\_DIGIT> | `_')*

<binaryLiteral> ::= `0b' <BINARY\_DIGIT> (<BINARY\_DIGIT> | `_')*

<HEX\_DIGIT> ::= <DIGIT>
  \alt `a' .. `f'
  \alt `A' .. `F'

<BINARY\_DIGIT> ::= `0' | `1'
\end{grammar}

\subsection{Floating Point Numbers}
\LMLabel{floatingLiterals}

\LMHash{}%
A \IndexCustom{floating literal}{literal!floating} is a decimal numeral
with a fractional part, an exponent, or both.

\begin{grammar}
<floatingLiteral> ::= <decimalLiteral> `.' <decimalLiteral> <EXPONENT>?
  \alt <decimalLiteral> `.' <EXPONENT>?
  \alt `.' <decimalLiteral> <EXPONENT>?
  \alt <decimalLiteral> <EXPONENT>

<EXPONENT> ::= (`e' | `E') (`+' | `-')? <decimalLiteral>
\end{grammar}

\subsection{Characters}
\LMLabel{characterLiterals}

\LMHash{}%
A \IndexCustom{character literal}{literal!character} is a single character
or escape sequence enclosed in single quotes.

\begin{grammar}
<characterLiteral> ::= `\sq' <characterContent> `\sq'

<characterContent> ::= <CHARACTER>
  \alt <escapeSequence>
\end{grammar}

\subsection{Strings}
\LMLabel{stringLiterals}

\LMHash{}%
A \Index{string literal} is a normal string, a multiline string,
or a raw string.
Normal and multiline strings support escape sequences
and interpolation of the form \code{\{expression\}}.
Doubled braces denote literal brace characters.
In a raw string no escapes or interpolations are processed.

\begin{grammar}
<stringLiteral> ::= <normalString>
  \alt <multilineString>
  \alt <rawString>

<normalString> ::= `"' (<stringCharacter> | <escapeSequence>
  \gnewline{} | <interpolation>)* `"'

<multilineString> ::= `"""' (<multilineCharacter> | <escapeSequence>
  \gnewline{} | <interpolation>)* `"""'

<rawString> ::= `r' `"' <RAW\_STRING\_CHARACTER>* `"'

<escapeSequence> ::= `\\' `n'
  \alt `\\' `r'
  \alt `\\' `t'
  \alt `\\' `b'
  \alt `\\' `f'
  \alt `\\' `0'
  \alt `\\' `\\'
  \alt `\\' `"'
  \alt `\\' `\sq'
  \alt `\\' `x' <HEX\_DIGIT> <HEX\_DIGIT>
  \alt `\\' `u' <HEX\_DIGIT> <HEX\_DIGIT> <HEX\_DIGIT> <HEX\_DIGIT>

<interpolation> ::= `{' <expression> `}'

<stringCharacter> ::= <STRING\_CHARACTER>

<multilineCharacter> ::= <MULTILINE\_CHARACTER>

<escapedBrace> ::= `{{' | `}}'
\end{grammar}

\commentary{%
  The EBNF escape rule lists \code{"\\n"} through
  \code{"\\u" + 4 hex digits}; the single-quote escape is rendered above as
  a backslash terminal followed by \code{\sq}, because a raw quote would
  terminate the terminal.%
}

\section{Operators}
\LMLabel{operatorsLexicon}

\LMHash{}%
The table below lists every operator token of Hella.
Their meaning in expression positions is specified in
\ref{expressions}; user-defined overloads are specified in
\ref{operatorDeclarations}.

\begin{grammar}
<operator> ::= `+'
  \alt `-'
  \alt `*'
  \alt `/'
  \alt `\%'
  \alt `<'
  \alt `<='
  \alt `>'
  \alt `>='
  \alt \IS{} | `is' `not'
  \alt \AND{} | \OR{} | \NOT{}
  \alt `\&'
  \alt `|'
  \alt `^'
  \alt `~'
  \alt `\ltlt'
  \alt `\gtgt'
  \alt `='
  \alt `+=' | `-=' | `*=' | `/=' | `\%='
  \alt `\&=' | `|=' | `^=' | `\ltlt=' | `\gtgt='
  \alt `++' | `--'
  \alt `?'
  \alt `:'
  \alt `??'
  \alt `?.'
  \alt `..'
  \alt `..='
\end{grammar}

\section{Types}
\LMLabel{types}

\LMHash{}%
A \Index{type} is a primary type followed by zero or more modifiers.
The modifiers are nullability (\code{?}), pointer (\code{*})
and slice/array view (\code{[]} ).

\begin{grammar}
<type> ::= <primaryType> <typeModifier>*

<primaryType> ::= <primitiveType>
  \alt <namedType>
  \alt <genericType>
  \alt <functionType>
  \alt <tupleType>
  \alt <anyType>
  \alt <fixedArrayType>
  \alt <vectorType>
  \alt <mapType>
  \alt <ownType>
  \alt <taskType>
  \alt <parenthesizedType>

<primitiveType> ::= \INT{} | \BOOL{} | \FLOAT{} | \DOUBLE{}
  \alt \code{char} | \STRING{} | \VOID{}

<anyType> ::= \ANY{}

<namedType> ::= <qualifiedName>

<qualifiedName> ::= <identifier> (`::' <identifier>)*

<genericType> ::= <qualifiedName> `<' <typeArguments> `>'

<typeArguments> ::= <type> (`,' <type>)*

<functionType> ::= \FUNCTION{} `<' <type>? `(' <typeList>? `)' `>'

<typeList> ::= <type> (`,' <type>)*

<tupleType> ::= `'(' (<type> (`,' <type>)* `,'?)? `')'

<parenthesizedType> ::= `'(' <type> `')'

<typeModifier> ::= `?'
  \alt `*'
  \alt `[' `]'
\end{grammar}

\LMHash{}%
\code{arr} is a compiler keyword for fixed-size arrays (element type
precedes it). Inferred size takes the initializer length; explicit
size is zero-initialized and fixed for the array's lifetime.

\begin{grammar}
<fixedArrayType> ::= <type> \ARR{} (`[' <integerLiteral> `]')?
\end{grammar}

\LMHash{}%
\code{vec} is a compiler keyword for owning, growable vectors (element
type precedes it). Empty vectors come from the \code{vec[]} expression
with an undetermined element type until the first \code{push}.

\begin{grammar}
<vectorType> ::= <type> \VEC{}
\end{grammar}

\LMHash{}%
There is no \code{map} keyword. Map types use \code{KEY\_TYPE:VALUE\_TYPE}.
The \code{:} operator is distinct from \code{::} (qualified names) and
from the \code{where} constraint colon.
\code{string:int[]} is a map to an array.

\begin{grammar}
<mapType> ::= <type> `:' <type>
\end{grammar}

\LMHash{}%
An owning heap pointer is written \code{own T}: it moves on assignment,
argument passing and return, \code{null} never flows into it,
\code{delete} destroys it early and poisons the slot, and scope exit
destroys any remaining live \code{own} slot. In phase 1 \code{own} never
nests inside other types and \code{own} fields and globals are rejected.

\begin{grammar}
<ownType> ::= \OWN{} <type>
\end{grammar}

\LMHash{}%
A pending asynchronous computation is written \code{task<T>}: it is the
value produced by calling an \code{async} function or by \code{spawn}
(\ref{asyncFunctions}). \code{await} consumes a \code{task<T>} handle
exactly once and yields \code{T} (\code{task<void>} awaits to no
value). \code{task} is not a user-declarable generic type: only the
single-type-argument form is valid.

\begin{grammar}
<taskType> ::= \code{task} `<' <type> `>'
\end{grammar}

\commentary{%
  The EBNF grammar is ambiguous between a parenthesized type and a
  one-element tuple type. As in PEGs, alternation gives priority to the
  left, so \synt{functionType} is tried before \synt{tupleType}, and
  \synt{mapType} binds through both operand types.%
}

\section{Generic Parameters and Constraints}
\LMLabel{genericParameters}

\LMHash{}%
Generic declarations may declare type parameters with optional bounds,
and functions may declare an additional \code{where} clause with
further constraints written \code{T: Bound1, Bound2}.

\begin{grammar}
<genericParameters> ::= `<' <genericParameter> (`,' <genericParameter>)* `>'

<genericParameter> ::= <identifier> (`:' <constraintList>)?

<constraintList> ::= <type> (`,' <type>)*

<whereClause> ::= \WHERE{} <whereConstraint> <whereConstraint>*

<whereConstraint> ::= <type> `:' <constraintList>
\end{grammar}

\section{Expressions}
\LMLabel{expressions}

\LMHash{}%
An \Index{expression} is an assignment expression.
Operator precedence rises from assignment (lowest) through conditional,
null-coalescing, logical (\code{or}, \code{and}), bitwise
(\code{|}, \code{\^}, \code{\&}), equality (\code{is}, \code{is not}),
relational, shift, additive and multiplicative (highest),
ending in unary and postfix forms.

\begin{grammar}
<expression> ::= <assignmentExpression>

<assignmentExpression> ::= <conditionalExpression>
  \alt <unaryExpression> <assignmentOperator> <assignmentExpression>

<assignmentOperator> ::= `='
  \alt `+=' | `-=' | `*=' | `/=' | `\%='
  \alt `\&=' | `|=' | `^=' | `\ltlt=' | `\gtgt='

<conditionalExpression> ::= <nullCoalescingExpression>
  \alt <nullCoalescingExpression> `?' <expression> `:' <conditionalExpression>

<nullCoalescingExpression> ::= <logicalOrExpression>
  \alt <logicalOrExpression> `??' <nullCoalescingExpression>

<logicalOrExpression> ::= <logicalAndExpression>
  \alt <logicalOrExpression> \OR{} <logicalAndExpression>

<logicalAndExpression> ::= <bitwiseOrExpression>
  \alt <logicalAndExpression> \AND{} <bitwiseOrExpression>

<bitwiseOrExpression> ::= <bitwiseXorExpression>
  \alt <bitwiseOrExpression> `|' <bitwiseXorExpression>

<bitwiseXorExpression> ::= <bitwiseAndExpression>
  \alt <bitwiseXorExpression> `^' <bitwiseAndExpression>

<bitwiseAndExpression> ::= <equalityExpression>
  \alt <bitwiseAndExpression> `\&' <equalityExpression>

<equalityExpression> ::= <relationalExpression>
  \alt <equalityExpression> <equalityOperator> <relationalExpression>

<equalityOperator> ::= \IS{} | `is' `not'

<relationalExpression> ::= <shiftExpression>
  \alt <relationalExpression> <relationalOperator> <shiftExpression>

<relationalOperator> ::= `<'
  \alt `<='
  \alt `>'
  \alt `>='

<shiftExpression> ::= <additiveExpression>
  \alt <shiftExpression> <shiftOperator> <additiveExpression>

<shiftOperator> ::= `\ltlt'
  \alt `\gtgt'

<additiveExpression> ::= <multiplicativeExpression>
  \alt <additiveExpression> <additiveOperator> <multiplicativeExpression>

<additiveOperator> ::= `+'
  \alt `-'

<multiplicativeExpression> ::= <unaryExpression>
  \alt <multiplicativeExpression> <multiplicativeOperator> <unaryExpression>

<multiplicativeOperator> ::= `*'
  \alt `/'
  \alt `\%'

<unaryExpression> ::= <unaryOperator> <unaryExpression>
  \alt <awaitExpression>
  \alt <spawnExpression>
  \alt <postfixExpression>

<awaitExpression> ::= \AWAIT{} <unaryExpression>

<spawnExpression> ::= \SPAWN{} <unaryExpression>

<unaryOperator> ::= `++' | `--' | `+' | `-' | `~' | \NOT{}
\end{grammar}

\LMHash{}%
Prefix \code{await} and \code{spawn} bind at the unary level (like
\code{not}), so \code{await a + b} requires parentheses around the
addition. Both are only legal inside an \code{async} function body
(\ref{asyncFunctions}): \code{await} consumes a \code{task<T>} handle
exactly once, \code{spawn} additionally requires an enclosing
\code{scope} block.

\subsection{Postfix and Primary Expressions}
\LMLabel{postfixExpressions}

\LMHash{}%
Postfix operations (calls, indexing, member access, nullable member
access and post-increment/decrement) bind tightest after primary
forms. A call may pass positional arguments, named arguments
(\code{name: expression}), and \code{out}/\code{ref} arguments.
Indexing selects an element or a range (\code{..}, \code{..=}).

\begin{grammar}
<postfixExpression> ::= <primaryExpression> <postfixOperation>*

<postfixOperation> ::= <callExpression>
  \alt <indexExpression>
  \alt <memberAccess>
  \alt <nullableMemberAccess>
  \alt <postfixIncrement>
  \alt <postfixDecrement>

<callExpression> ::= `(' <argumentList>? `)'

<argumentList> ::= <argument> (`,' <argument>)*

<argument> ::= <expression>
  \alt <namedArgument>
  \alt \OUT{} <type>? <identifier>
  \alt \REF{} <expression>

<namedArgument> ::= <identifier> `:' <expression>

<indexExpression> ::= `[' <indexOrRange> `]'

<indexOrRange> ::= <expression>
  \alt <rangeExpression>
  \alt <expression>? <rangeOperator> <expression>?

<rangeExpression> ::= <expression> <rangeOperator> <expression>

<rangeOperator> ::= `..' | `..='

<memberAccess> ::= `.' <identifier>

<nullableMemberAccess> ::= `?.' <identifier>

<postfixIncrement> ::= `++'

<postfixDecrement> ::= `--'
\end{grammar}

\begin{grammar}
<primaryExpression> ::= <literal>
  \alt <identifierExpression>
  \alt <qualifiedExpression>
  \alt \THIS{} | \SUPER{}
  \alt <closureExpression>
  \alt <matchExpression>
  \alt <structExpression>
  \alt <enumVariantExpression>
  \alt <tupleExpression>
  \alt <arrayLiteral>
  \alt <emptyVector>
  \alt <mapLiteral>
  \alt <newExpression>
  \alt <parenthesizedExpression>

<identifierExpression> ::= <identifier>

<qualifiedExpression> ::= <qualifiedName>

<literal> ::= <integerLiteral>
  \alt <floatingLiteral>
  \alt <characterLiteral>
  \alt <stringLiteral>
  \alt `true' | `false' | \NULL{}

<parenthesizedExpression> ::= `(' <expression> `)'

<tupleExpression> ::= `(' (<expression> (`,' <expression>)* `,'?)? `)'
\end{grammar}

\subsection{Collection Literals}
\LMLabel{collectionLiterals}

\LMHash{}%
An array literal is the initializer for a \code{TYPE arr} fixed-size
array (size inferred from the element count). All elements must be
compatible with the element type.

\begin{grammar}
<arrayLiteral> ::= `[' (<expression> (`,' <expression>)*)? `]'
\end{grammar}

\LMHash{}%
The empty vector \code{vec[]} has no element type until the first
\code{push} establishes it permanently (for example,
\code{any xs = vec[]} takes its element type from the first push).

\begin{grammar}
<emptyVector> ::= \VEC{} `[' `]'
\end{grammar}

\LMHash{}%
A map literal holds \code{key: value} entries inside a \code{has}
block for a \code{KEY:VALUE} type. Entries may be written multiline
or compact (\code{has "a": 1, "b": 2 end}); the declared type may
be omitted for \code{any} (inferred from the first entry). Keys must
be \code{string} or int types; both sides are type-checked.

\begin{grammar}
<mapLiteral> ::= <type>? \HAS{} <mapEntry>* \END{}

<mapEntry> ::= <expression> `:' <expression> `,'? <statementTerminator>
\end{grammar}

\section{Closures}
\LMLabel{closures}

\LMHash{}%
A closure abstracts over a parameter list and a body.
The body is either a single expression after \code{=>}
or a block.

\begin{grammar}
<closureExpression> ::= `|' <closureParameters>? `|' <closureBody>

<closureParameters> ::= <closureParameter> (`,' <closureParameter>)*

<closureParameter> ::= <identifier>
  \alt <type> <identifier>

<closureBody> ::= `=>' <expression>
  \alt <block>
\end{grammar}

\section{Match}
\LMLabel{matchExpressions}

\LMHash{}%
A match expression scrutinizes a value against a sequence of arms.
Each arm pairs a pattern with an optional \code{if} guard and a body
which is either an expression or a block.
The match statement (\ref{matchStatement}) shares the same arm grammar.

\begin{grammar}
<matchExpression> ::= \MATCH{} <expression> \DO{} <matchArm>* \END{}

<matchArm> ::= <pattern> <guard>? `->' <matchArmBody>

<matchArmBody> ::= <expression>
  \alt <block>

<guard> ::= \IF{} <expression>
\end{grammar}

\LMHash{}%
A pattern is an alternative chain of primary patterns separated by
\code{|} or its keyword alias \code{or} (\code{1 | 2} is equivalent
to \code{1 or 2}).

\begin{grammar}
<pattern> ::= <patternAlternative>

<patternAlternative> ::= <patternPrimary> ((`|' | \OR{}) <patternPrimary>)*

<patternPrimary> ::= `_'
  \alt <literal>
  \alt <identifier>
  \alt <enumPattern>
  \alt <tuplePattern>

<enumPattern> ::= `.' <identifier> (`(' <patternList>? `)')?
  \alt <qualifiedName> `.' <identifier> (`(' <patternList>? `)')?
  \alt <identifier> `(' <patternList>? `)'

<patternList> ::= <pattern> (`,' <pattern>)*

<tuplePattern> ::= `(' <pattern> (`,' <pattern>)* `,'? `)'
\end{grammar}

\section{Struct Construction}
\LMLabel{structConstruction}

\LMHash{}%
A struct expression initializes a struct (or constructor-less class)
with a \code{has} block of field assignments.
The leading type may be omitted when the expression appears as the
initializer of a variable declaration whose declared type is a struct
or constructor-less class (not \code{any}). In
\code{User u2 = has name = "bbb" end} the second \code{has} infers
\code{User}.

\begin{grammar}
<structExpression> ::= <type>? \HAS{} <structFieldInitializer>* \END{}

<structFieldInitializer> ::= <identifier> `=' <expression> <statementTerminator>

<newExpression> ::= \NEW{} <type> `(' <argumentList>? `)'
\end{grammar}

\LMHash{}%
Heap construction \code{new Type(args)} allocates a fresh heap object,
runs the constructor on it, and yields an owning value.
The result must be stored into an \code{own T} slot; a bare
\code{new} in an expression-statement position leaks and is rejected.

\section{Enum Variant Construction}
\LMLabel{enumVariants}

\LMHash{}%
An enum variant is named with a leading dot, optionally applied to
arguments, optionally qualified by the enum type name.

\begin{grammar}
<enumVariantExpression> ::= `.' <identifier>
  \alt `.' <identifier> `(' <argumentList>? `)'
  \alt <qualifiedName> `.' <identifier>
  \alt <qualifiedName> `.' <identifier> `(' <argumentList>? `)'
\end{grammar}

\section{Statements}
\LMLabel{statements}

\LMHash{}%
A statement is an expression statement, a declaration, a control-flow
statement, an assertion, a heap-destruction statement, or a block.
A statement terminator is a newline or a semicolon.

\begin{grammar}
<statement> ::= <expressionStatement>
  \alt <variableDeclaration>
  \alt <constantDeclaration>
  \alt <destructuringDeclaration>
  \alt <destructuringAssignment>
  \alt <ifStatement>
  \alt <matchStatement>
  \alt <whileStatement>
  \alt <loopStatement>
  \alt <forStatement>
  \alt <breakStatement>
  \alt <continueStatement>
  \alt <returnStatement>
  \alt <deferStatement>
  \alt <deleteStatement>
  \alt <asyncScopeStatement>
  \alt <yieldStatement>
  \alt <assertStatement>
  \alt <block>

<expressionStatement> ::= <expression> <statementTerminator>

<statementTerminator> ::= <NEWLINE>
  \alt `;'

<block> ::= \DO{} <statement>* \END{}
\end{grammar}

\section{Variable Declarations}
\LMLabel{variableDeclarations}

\LMHash{}%
A variable declaration optionally carries a visibility modifier
(\ref{modules}), a type, a name, and an optional initializer.
A constant declaration uses \code{const} with an optional type.
Destructuring binds several names (or discards with \code{\_})
from a single right-hand side.

\begin{grammar}
<variableDeclaration> ::= <visibilityModifier>? <type> <identifier>
  \gnewline{} (`=' <expression>)? <statementTerminator>

<constantDeclaration> ::= <visibilityModifier>? \CONST{} <type>? <identifier>
  \gnewline{} `=' <expression> <statementTerminator>

<destructuringDeclaration> ::= <destructuringTarget>
  \gnewline{} `=' <expression> <statementTerminator>

<destructuringTarget> ::= <identifier> (`,' (<identifier> | `_'))*

<destructuringAssignment> ::= <destructuringTarget>
  \gnewline{} `=' <expression> <statementTerminator>
\end{grammar}

\section{Conditionals}
\LMLabel{conditionals}

\LMHash{}%
An if statement evaluates its condition and executes the matching block.
Parentheses around conditions are optional, and any number of
\code{else if} branches may precede a final \code{else} block.

\begin{grammar}
<ifStatement> ::= \IF{} (`(')? <expression> (`)')? <block>
  \gnewline{} (\ELSE{} \IF{} (`(')? <expression> (`)')? <block>)*
  \gnewline{} (\ELSE{} <block>)?
\end{grammar}

\section{Match Statement}
\LMLabel{matchStatement}

\LMHash{}%
A match statement has the same form as a match expression
(\ref{matchExpressions}) but executes arm bodies as statements.

\begin{grammar}
<matchStatement> ::= \MATCH{} <expression> \DO{} <matchArm>* \END{}
\end{grammar}

\section{Loops}
\LMLabel{loops}

\LMHash{}%
A while loop takes an optionally parenthesized condition.
A loop is optionally labeled and executes its block repeatedly.
A for loop iterates over an expression: \code{for x in arr} binds
elements and \code{for x, i in arr} also binds the index. Over maps,
\code{for k in m} binds keys and \code{for k, v in m} binds
key/value pairs.

\begin{grammar}
<whileStatement> ::= \WHILE{} (`(')? <expression> (`)')? <block>

<loopStatement> ::= <loopLabel>? \LOOP{} <block>

<forStatement> ::= <loopLabel>? \FOR{} <identifier> (`,' <identifier>)?
  \gnewline{} \IN{} <expression> <block>

<loopLabel> ::= <identifier> `:'

<breakStatement> ::= \BREAK{} <identifier>? <statementTerminator>

<continueStatement> ::= \CONTINUE{} <identifier>? <statementTerminator>
\end{grammar}

\section{Return}
\LMLabel{returnStatement}

\LMHash{}%
A return statement optionally carries a value.

\begin{grammar}
<returnStatement> ::= \RETURN{} <expression>? <statementTerminator>
\end{grammar}

\section{Defer}
\LMLabel{deferStatement}

\LMHash{}%
A defer statement runs its expression or block when the enclosing
scope exits.

\begin{grammar}
<deferStatement> ::= \DEFER{} (<expression> <statementTerminator> | <block>)

<deleteStatement> ::= \DELETE{} <expression> <statementTerminator>
\end{grammar}

\LMHash{}%
A delete statement destroys the owning heap value denoted by its operand
early and poisons the slot. The operand must be an \code{own T} slot
(an identifier); \code{delete this} and non-identifier operands are
rejected.

\LMHash{}%
An async scope statement opens a structured-concurrency scope: tasks
declared inside it must be awaited before \code{end}, and scope exit
joins the remaining children (\ref{asyncFunctions}). It is only legal
directly inside an \code{async} function body; nesting is allowed.

\begin{grammar}
<asyncScopeStatement> ::= \SCOPE{} <block>
\end{grammar}

\LMHash{}%
A yield statement is a cooperative checkpoint inside an \code{async}
body (\ref{asyncFunctions}). It is an \Error{compile-time error} to use
\code{scope} or \code{yield} outside an \code{async} function body.

\commentary{%
  The source EBNF lists these as standalone \code{async-scope-statement}
  and \code{yield-statement} productions without referencing them from
  \code{statement}; they are presented here as statement alternatives,
  which is how the implementation accepts them.%
}

\begin{grammar}
<yieldStatement> ::= \YIELD{} <statementTerminator>?
\end{grammar}

\section{Assertions}
\LMLabel{assertions}

\LMHash{}%
An assertion checks its condition, with an optional message.
\code{debug\_assert} is only checked in debug builds.

\begin{grammar}
<assertStatement> ::= \ASSERT{} <expression> (`,' <expression>)?
  \gnewline{} <statementTerminator>
  \alt \DEBUGASSERT{} <expression> (`,' <expression>)?
  \gnewline{} <statementTerminator>
\end{grammar}

\section{Functions}
\LMLabel{functions}

\LMHash{}%
A function declaration carries an optional visibility modifier,
optional \code{static}, \code{sealed} and \code{override} modifiers,
a return type, a name, optional generic parameters, a parameter list,
an optional \code{where} clause, an optional \code{async} marker
(\ref{asyncFunctions}), and a body which is either a block
or an \code{initialize} stub.

\begin{grammar}
<functionDeclaration> ::= <visibilityModifier>? \STATIC{}? \SEALED{}?
  \gnewline{} \OVERRIDE{}? <type> <identifier> <genericParameters>?
  \gnewline{} `(' <parameterList>? `)' <whereClause>? \ASYNC{}? <functionBody>

<functionBody> ::= <block>
  \alt \INITIALIZE{} <statementTerminator>

<parameterList> ::= <parameter> (`,' <parameter>)*

<parameter> ::= `...' ? <parameterMode>? <type>? <identifier>
  \gnewline{} (`=' <expression>)?

<parameterMode> ::= \REF{}
  \alt \OUT{}

<variadicMarker> ::= `...'
\end{grammar}

\commentary{%
  \code{...} marks variadic parameters: \code{...int vda} yields
  \code{vda: int[]}, \code{...T vda} yields \code{vda: T[]} with a
  \code{where T: Trait} bound allowed, and bare \code{... vda} (no type)
  derives \code{vda: prev\_type[]} from the previous type and must be last;
  \code{...} may appear anywhere when an explicit type is present, e.g.
  \code{log(string fmt, ...string vda, bool cond)}.%
}

\subsection{Async Functions}
\LMLabel{asyncFunctions}

\LMHash{}%
A function declaration may carry the \code{async} marker between its
signature and its body: \code{ReturnType name(params) async do ... end}.
An \code{async} body runs as its own task. Calling an \code{async}
function from an \code{async} body yields a \code{task<Ret>} handle;
\code{spawn} is the explicit form and additionally requires an
enclosing \code{scope} block.

\LMHash{}%
Task handles are single-consumer and scope-bound: \code{await} consumes
a handle exactly once (a second \code{await} of the same handle is an
\Error{compile-time error}), and every task declared inside a
\code{scope} block must be awaited before \code{end} (scope exit is the
join point; an escaping task is an \Error{compile-time error}).
Dropping a handle never detaches work. An unconsumed non-void task
result at scope exit is diagnosed.

\LMHash{}%
It is an \Error{compile-time error} to call an \code{async} function
from a synchronous body, and to use \code{await}, \code{spawn},
\code{yield} or \code{scope} outside an \code{async} body.
Arguments to a spawned child are transferred by value. An \code{async}
\code{main} runs as the root task (\ref{programEntry}).

\commentary{%
  In version 1, \code{async} is supported on free functions only:
  methods, trait signatures and function types remain synchronous, and
  the async runtime is linked only when reachable async code requires
  it, so synchronous programs gain no scheduler symbols or
  dependencies. The standard library offers timer/yield helpers
  (\code{yieldNow()}, \code{sleepMs(int)}, \code{cancelled()});
  genuinely asynchronous I/O, channels, async closures and detached
  tasks are future work.%
}

\section{Classes}
\LMLabel{classes}

\LMHash{}%
\code{open} is only for \code{class} (it enables \code{extend});
methods use \code{static}/\code{sealed}/\code{override} but never
\code{open}: \code{open int bar()} is rejected. Constructors and
\code{get}/\code{set} accessors are \code{public} by default; other
members are \code{private} unless marked \code{public}.

\begin{grammar}
<classDeclaration> ::= <visibilityModifier>? \OPEN{}? \CLASS{} <identifier>
  \gnewline{} <genericParameters>? (\EXTENDS{} <type>)?
  \gnewline{} (\IMPLEMENTS{} <traitList>)? \HAS{} <classMember>* \END{}

<traitList> ::= <type> (`,' <type>)*

<classMember> ::= <fieldDeclaration>
  \alt <functionDeclaration>
  \alt <constructorDeclaration>
  \alt <destructorDeclaration>
  \alt <propertyDeclaration>
  \alt <operatorDeclaration>
  \alt <conversionDeclaration>
  \alt <constantDeclaration>

<fieldDeclaration> ::= <visibilityModifier>? <type> <identifier>
  \gnewline{} (`=' <expression>)? <statementTerminator>

<constructorDeclaration> ::= <visibilityModifier>? <identifier>
  \gnewline{} `(' <parameterList>? `)' \INITIALIZE{} <block>?

<destructorDeclaration> ::= <visibilityModifier>? `~' <identifier>
  \gnewline{} `(' `)' <block>
\end{grammar}

\section{Properties}
\LMLabel{properties}

\LMHash{}%
\code{get}/\code{set} are \code{public} by default (overriding the
class \code{private} default). Two property declarations with the same
name in one class merge if they supply complementary accessors:
\code{int x get ... end} plus \code{int x set(int v) ... end} becomes
one property.

\begin{grammar}
<propertyDeclaration> ::= <visibilityModifier>? <type>? <identifier>
  \gnewline{} <propertyGetter> <propertySetter>?
  \alt <visibilityModifier>? <type>? <identifier>
  \gnewline{} <propertySetter> <propertyGetter>?

<propertyGetter> ::= \GET{} <block>

<propertySetter> ::= \SET{} `(' <parameter> `)' <block>
\end{grammar}

\section{Structs}
\LMLabel{structs}

\LMHash{}%
A struct declaration groups typed fields in a \code{has} block.

\begin{grammar}
<structDeclaration> ::= <visibilityModifier>? \STRUCT{} <identifier>
  \gnewline{} <genericParameters>? \HAS{} <structMember>* \END{}

<structMember> ::= <visibilityModifier>? <type> <identifier>
  \gnewline{} (`=' <expression>)? <statementTerminator>
\end{grammar}

\section{Traits}
\LMLabel{traits}

\LMHash{}%
A trait declaration groups function signatures in a \code{has} block.

\begin{grammar}
<traitDeclaration> ::= <visibilityModifier>? \TRAIT{} <identifier>
  \gnewline{} <genericParameters>? \HAS{} <traitMember>* \END{}

<traitMember> ::= <functionSignature>

<functionSignature> ::= \SEALED{}? <type> <identifier>
  \gnewline{} <genericParameters>? `(' <parameterList>? `)'
  \gnewline{} <whereClause>? <statementTerminator>
\end{grammar}

\section{Operator Declarations}
\LMLabel{operatorDeclarations}

\LMHash{}%
An operator declaration overloads one of the operator symbols below for
a type. Indexing overloads use \code{[]} (no \code{[]=} symbol here;
assignment forms live in expression grammar).

\begin{grammar}
<operatorDeclaration> ::= <visibilityModifier>? \STATIC{}? \OPERATOR{}
  \gnewline{} <operatorSymbol> `(' <parameterList>? `)'
  \gnewline{} <whereClause>? <block>

<operatorSymbol> ::= `+'
  \alt `-'
  \alt `*'
  \alt `/'
  \alt `\%'
  \alt `<'
  \alt `<='
  \alt `>'
  \alt `>='
  \alt \IS{} | `is' `not'
  \alt `\&'
  \alt `|'
  \alt `^'
  \alt `~'
  \alt `\ltlt'
  \alt `\gtgt'
  \alt `='
  \alt `+=' | `-=' | `*=' | `/=' | `\%='
  \alt `\&=' | `|=' | `^=' | `\ltlt=' | `\gtgt='
  \alt `++' | `--'
  \alt `[' `]'
\end{grammar}

\section{Conversions}
\LMLabel{conversions}

\LMHash{}%
A conversion declaration defines a (possibly \code{explicit})
conversion from one type to another.

\begin{grammar}
<conversionDeclaration> ::= <visibilityModifier>? \EXPLICIT{}? \CONVERT{}
  \gnewline{} <type> \TO{} <type> <block>
\end{grammar}

\section{Enums}
\LMLabel{enums}

\LMHash{}%
An enum declaration groups variants in a \code{has} block.
Each variant optionally carries a discriminant value and an optional
parameter list for associated data.

\begin{grammar}
<enumDeclaration> ::= <visibilityModifier>? \ENUM{} <identifier>
  \gnewline{} <genericParameters>? \HAS{} <enumVariant>* \END{}

<enumVariant> ::= <identifier> (`=' <expression>)?
  \gnewline{} (`(' <parameterList>? `)')? <statementTerminator>
\end{grammar}

\section{Type Aliases}
\LMLabel{typeAliases}

\LMHash{}%
A type alias introduces a name for a type, with optional generic
parameters.

\begin{grammar}
<typedefDeclaration> ::= <visibilityModifier>? \TYPEDEF{}
  \gnewline{} <genericParameters>? <identifier> `=' <type> <statementTerminator>
\end{grammar}

\section{Distinct Types}
\LMLabel{distinctTypes}

\LMHash{}%
A distinct type introduces a new nominal type with the representation
of its aliased type.

\begin{grammar}
<distinctDeclaration> ::= <visibilityModifier>? \DISTINCT{} <identifier>
  \gnewline{} <genericParameters>? `=' <type> <statementTerminator>
\end{grammar}

\section{Extensions}
\LMLabel{extensions}

\LMHash{}%
Extension members are public by default: an absent visibility modifier
normalizes to \code{public} when merged into the target type.
An explicit \code{private} is kept and still enforced.

\begin{grammar}
<extensionDeclaration> ::= \EXTEND{} <type> \DO{}
  \gnewline{} <extensionMember>* \END{}

<extensionMember> ::= <fieldDeclaration>
  \alt <functionDeclaration>
  \alt <operatorDeclaration>
  \alt <propertyDeclaration>
  \alt <conversionDeclaration>
\end{grammar}

\section{Modules}
\LMLabel{modules}

\LMHash{}%
An import declaration names a module path as a qualified name, with an
optional brace-enclosed list of imported members.
Visibility is either \code{public} or \code{private}.

\begin{grammar}
<importDeclaration> ::= \IMPORT{} <qualifiedName>
  \gnewline{} (`::' `{' <importList> `}')? <statementTerminator>

<importList> ::= <identifier> (`,' <identifier>)*

<visibilityModifier> ::= \PUBLIC{}
  \alt \PRIVATE{}
\end{grammar}

\section{Module Initialization}
\LMLabel{moduleInitialization}

\LMHash{}%
An init block runs once when its module is initialized.

\begin{grammar}
<initBlock> ::= \INIT{} <block>
\end{grammar}

\section{Attributes}
\LMLabel{attributes}

\LMHash{}%
Attributes (annotations) attach metadata to a top-level declaration.
Each argument is an expression, a \code{name: expression} pair,
or a \code{name = expression} pair.

\begin{grammar}
<attribute> ::= `@' <identifier> (`(' <attributeArguments>? `)')?

<attributeArguments> ::= <attributeArgument> (`,' <attributeArgument>)*

<attributeArgument> ::= <expression>
  \alt <identifier> `:' <expression>
  \alt <identifier> `=' <expression>

<attributedDeclaration> ::= <attribute>* <topLevelDeclaration>
\end{grammar}

\subsection{Conditional Compilation}
\LMLabel{conditionalCompilation}

\LMHash{}%
A \code{@cfg(...)} attribute keeps its item only when its condition
holds. Evaluation happens during import expansion, before semantic
analysis, so an absent item is invisible to name resolution, checking
and code generation, and mutually exclusive declarations of the same
name never collide.

\LMHash{}%
The version 1 conditions are: \code{os = "<name>"} for the host OS
(\code{macos} \code{|} \code{linux} \code{|} \code{windows}),
\code{target = "<triple>"} for the exact host triple,
\code{debug} (true only for a debug link profile),
\code{debug = false} (true only for release, check and LSP analysis),
and the \code{unix} \code{|} \code{windows} OS-family shorthands.
\code{or} (or \code{,}) joins conditions disjunctively: any true
condition keeps the item. Stacked attributes are conjunctive: every
\code{@cfg} must hold. There is no \code{and} or \code{not} in version
1. Unknown conditions evaluate to false.

\begin{grammar}
<cfgCondition> ::= <identifier>
  \alt <identifier> `=' (<stringLiteral> | `true' | `false')

<cfgArguments> ::= <cfgCondition> ((`,' | \OR{}) <cfgCondition>)*
\end{grammar}

\commentary{%
  \code{check}, \code{lint} and the language server evaluate with
  \code{debug = false} so analysis never depends on the link profile.
  The canonical formatter rewrites \code{or} as a comma (both mean
  ``any of''); it never rewrites the author's \code{=} into \code{:}
  or vice versa.%
}

\section{Top-Level Declarations}
\LMLabel{topLevelDeclarations}

\LMHash{}%
A source file is built from the declarations below.
Attributes and module initialization blocks attach to (respectively
accompany) these forms.

\begin{grammar}
<topLevelDeclaration> ::= <functionDeclaration>
  \alt <classDeclaration>
  \alt <structDeclaration>
  \alt <traitDeclaration>
  \alt <enumDeclaration>
  \alt <typedefDeclaration>
  \alt <distinctDeclaration>
  \alt <extensionDeclaration>
  \alt <externDeclaration>
  \alt <importDeclaration>
  \alt <constantDeclaration>
  \alt <variableDeclaration>
  \alt <attributedDeclaration>
\end{grammar}

\commentary{%
  The source EBNF omits \code{extern-declaration} from
  \code{top-level-declaration}; it is included here because the
  implementation accepts bare \code{extern} blocks at top level.%
}

\section{Foreign Function Interface}
\LMLabel{ffi}

\LMHash{}%
An extern declaration binds foreign functions, structs, enums and
constants from a named foreign module.
Extern members are always public: an explicit \code{public} modifier
on a member is accepted and ignored, \code{private} is rejected.
No visibility is stored --- nothing gates on extern visibility.

\begin{grammar}
<externDeclaration> ::= <visibilityModifier>? \EXTERN{} <stringLiteral>
  \gnewline{} \FROM{} <stringLiteral> \DO{} <externMember>* \END{}

<externMember> ::= <externFunction>
  \alt <externStruct>
  \alt <externEnum>
  \alt <externConstant>

<externFunction> ::= <type> <identifier>
  \gnewline{} `(' <externParameterList>? `)' <statementTerminator>

<externParameterList> ::= (<externParameter> (`,' <externParameter>)*)?
  \gnewline{} (`,' `...')? | `...'

<externParameter> ::= <parameterMode>? `...'? <type>? <identifier>
  \alt `...'
\end{grammar}

\LMHash{}%
An extern parameter may carry a \code{ref} or \code{out} mode for C
functions that write through a pointer (for example,
\code{int rand\_s(ref int out)}): the call site passes the variable's
address, exactly like a Hella \code{ref} parameter, and the mode is
enforced against the call-site marker. A mode does not combine with
\code{...}.

\commentary{%
  Bare \code{...} covers C varargs like
  \code{printf(string fmt, ...)}; \code{...int vda} yields
  \code{vda: int[]}. The source EBNF does not yet list the
  \code{ref}/\code{out} parameter mode; it is included here because the
  implementation accepts and enforces it.%
}

\begin{grammar}
<externStruct> ::= \STRUCT{} <identifier>
  \gnewline{} (\HAS{} <externField>* \END{})?

<externField> ::= <type> <identifier> <statementTerminator>

<externEnum> ::= \ENUM{} <identifier> \HAS{} <enumVariant>* \END{}

<externConstant> ::= \CONST{} <type> <identifier> <statementTerminator>
\end{grammar}

\section{Program Entry}
\LMLabel{programEntry}

\LMHash{}%
A program starts in one of the two \code{main} forms below.

\begin{grammar}
<mainFunction> ::= \VOID{} `main' `(' `)' <block>
  \alt \INT{} `main' `(' \code{string} `[' `]' `args' `)' <block>
\end{grammar}

\commentary{%
  The second form receives its command-line arguments as
  \code{string[] args}; the EBNF terminal \code{"string[]"} denotes the
  three-character type suffix form, not a distinct keyword.
  Either form may carry the \code{async} marker (\ref{asyncFunctions}),
  in which case it runs as the root task.
  Independently of the return type, the implementation also accepts each
  form with an empty parameter list or with \code{string[] args}.%
}

\section{Special Names}
\LMLabel{specialNames}

\LMHash{}%
\code{this}, \code{super} and \code{Self} have special meaning and
cannot be used as ordinary identifiers in the positions where the
grammar expects them.

\begin{grammar}
<specialName> ::= \THIS{}
  \alt \SUPER{}
  \alt \SELF{}
\end{grammar}

\section{Identifier Reference}
\LMLabel{identifierReference}

\LMHash{}%
An \Index{identifier reference} is an occurrence of \synt{identifier}
that refers to a declaration in scope.
Reserved words (\ref{reservedWords}) cannot be used as identifiers.
The contextual words \code{get}, \code{set}, \code{has},
\code{end}, \code{convert}, \code{to}, \code{extend}, \code{from},
\code{initialize} and \code{Self} are only special in the syntactic
positions given by the grammar; elsewhere they are ordinary
identifiers.

\section{Reference}
\LMLabel{reference}

\subsection{Lexical Rules}
\LMLabel{lexicalRules}

\LMHash{}%
Hella source text is represented as a sequence of Unicode code points.
This sequence is first converted into a sequence of tokens
according to the lexical rules given in this specification.
At any point in the tokenization process,
the longest possible token is recognized.

\commentary{%
  This section collects the lexical productions (upper-case names)
  alongside the syntactic productions above, in the Reference chapter.%
}

\begin{grammar}
<CHARACTER> ::= \gtilde(`\sq' | `\\r' | `\\n')

<STRING\_CHARACTER> ::= \gtilde(`"' | `\\' | `\\r' | `\\n')

<MULTILINE\_CHARACTER> ::= \gtilde(`\\r' | `\\n')

<RAW\_STRING\_CHARACTER> ::= \gtilde(`"' | `\\r' | `\\n')

<ANY\_CHARACTER> ::= \gtilde(`\\r' | `\\n') | <LINE\_BREAK>

<LINE\_BREAK> ::= <LINE\_FEED> | <CARRIAGE\_RETURN\_LINE\_FEED>
  \alt <CARRIAGE\_RETURN>

<LINE\_FEED> ::= `\\n'

<CARRIAGE\_RETURN> ::= `\\r'

<CARRIAGE\_RETURN\_LINE\_FEED> ::= `\\r' `\\n'
\end{grammar}

\commentary{%
  The EBNF terminals \code{LF}, \code{CRLF}, \code{CR},
  \code{any-character} and \code{any-character-except-newline} are primitive
  lexical tokens. The single underscore in \code{identifier},
  \code{CHARACTER} string content and \code{RAW} strings is written
  \code{`_'} (the lexical rule for identifiers above uses \code{`_'} for the underscore).%
}

\printindex
\end{document}
